\subsection{Excluding proper index-five glue at rank eight}
\label{subsec:k12-r8-glue-five}

\begin{proposition}[K17-K12-R8-GLUE5]\label{prop:k12-r8-glue-five}
Let \(L\) be a rank-eight lattice whose complete minimum cone is
simplicial of dimension thirteen.  If \(L_0\) is the sublattice
generated by its minimum vectors, then \([L:L_0]\ne5\).
\end{proposition}

\begin{proof}
There are exactly thirteen minimum pairs.  By
Corollary~\ref{cor:k12-r8-index-five}, every basis selected from
the minimum vectors has index at most five in \(L\).
For an index-five frame, signs, permutation of the frame and a
change of quotient generator put the numerator in the form
\[
                         (1,1,1,1,1,2,2,2)/5.              \tag{G5.1}
\]
For clarity, this exact signed-coordinate statement follows from
the direct calculation in
\cite[\S9, case \(d=5\), p.~31; Proposition~9.1 and full proof,
p.~33]{MartinetIndex2001}.
The only possible coefficient distributions are
\((6,2),(5,3),(4,4)\); the first and last force sixteen
minimum pairs.  The same count is recorded in
\cite[complement to \S7, p.~4]{MartinetIndexErrata2014}.

Suppose \([L:L_0]=5\).  Every independent eight-element minimum
frame has index in \(L\) a positive multiple of five and at most
five.  It is therefore a \(\mathbb Z\)-basis of \(L_0\).
Fix one such frame, and let \(V\) be the \(8\)-by-\(13\) matrix
of representatives of the minimum pairs in its coordinates.
The matrix contains \(I_8\), and every maximal minor is zero or
\(\pm1\).  Any smaller nonzero square minor can be padded to
a maximal minor using complementary identity columns.
Thus \(V\) is totally unimodular.
Its rational column matroid is the same as its column matroid
over \(\mathbb F_2\), since every nonzero minor is \(\pm1\).
It has rank eight and nullity five.  It has no loops or
parallel elements, since the minimum pairs are nonzero and
distinct up to sign.

Fix one nonzero generator of \(L/L_0\), and represent it by
\(u/5\), where \(u\in L_0/5L_0\).
In every minimum basis \(B\) of \(L_0\), all eight coordinates
of \(u\) are nonzero and their product is a nonsquare in
\(\mathbb F_5\).  Indeed (G5.1) has product \(2^3\), a
nonsquare.  Signs are squares modulo five, permutations
preserve the product, and a change of quotient generator
multiplies the product by \(c^8\), which is a square.
The property therefore holds for the one fixed generator
in every basis.

Let \(v\) be a column outside a chosen basis, and let
\(S=\{j:v_j\ne0\}\).  Each \(v_j\) lies in \(\{0,1,-1\}\),
by the maximal-minor calculation.  For each \(i\in S\),
replacing the \(i\)th basis vector by \(v\) gives another
minimum basis.  If \(u\) has old coordinates \(a_j\), its
new coordinates are
\[
 a'_i=a_i/v_i,\qquad
 a'_j=a_j-v_j a_i/v_i\quad(j\ne i).
                                                               \tag{G5.2}
\]
For \(j\in S\), put \(r_j=a_j/v_j\in\mathbb F_5^\times\).
Nonvanishing of the new coordinates for every pivot implies
that these residues are pairwise distinct.  Hence \(|S|\le4\).
Let \(\chi_5\) denote the quadratic character of
\(\mathbb F_5^\times\).
Comparing the nonsquare products in the two bases and
canceling the unchanged coordinates gives
\[
 \chi_5\left(\prod_{j\in S}r_j\right)
 =
 \chi_5\left(r_i\prod_{\substack{j\in S\\j\ne i}}(r_j-r_i)\right)
 \qquad(i\in S).                                             \tag{G5.3}
\]
All sign factors disappear because \(-1\) is a square
modulo five.

No subset of size two or three satisfies (G5.3).
Multiplying all residues by \(c\ne0\) multiplies both
products by \(c^{|S|}\), preserving equality or inequality
of their characters.  The two multiplicative orbits of
two-element subsets are represented by \(\{1,4\}\) and
\(\{1,2\}\).  Pivoting at \(1\), the old and new products
are respectively \(4,3\) and \(2,1\), of opposite quadratic
characters.  The four three-element subsets form one
multiplicative orbit, represented by \(\{1,2,3\}\).
At pivot \(1\), its products are \(1,2\), again of opposite
characters.

Support one would make \(v\) a duplicate minimum pair;
support zero is impossible.  Hence every nonbasis column,
relative to every minimum basis, has support exactly four.
Every fundamental circuit consequently has five elements.
Any circuit is fundamental for a basis obtained by extending
that circuit with one element removed.  Therefore every
circuit of this binary matroid has size five.

Two distinct circuits cannot intersect.  Their symmetric
difference would be a nonzero binary cycle of positive even
size at most eight.  Every binary cycle is a disjoint union
of circuits: choose a minimal nonempty cycle in its support,
subtract it over \(\mathbb F_2\), and repeat.
If every circuit has size five, a nonzero cycle of even size
must contain at least two disjoint circuits and have at
least ten elements.  This is a contradiction.

All circuits are therefore pairwise disjoint.  Their
characteristic vectors are linearly independent by
disjointness, and span the binary kernel by the same
cycle decomposition.  Nullity five would require five
disjoint five-element circuits, hence at least twenty-five
columns.  There are only thirteen.  The contradiction
proves the proposition.
\end{proof}

The conclusion excludes proper ambient index five for the
remaining simplicial rank-eight cells.  It leaves
minimum-generated cells of maximal minimum-basis index five
and proper glue of the other allowed orders.  It does not
assert vanishing of their chains or homology.
